\paragraph{6. Central trace and period normalization.}
First use (E.1) to remove the unused tame primes. The central
specialization of (K8) is the corestriction of $z_D\otimes e_\chi$
from $\mathbf Q(\zeta_D)$ to $\mathbf Q$, with omitted support
$\operatorname{prime}(2Nd)$. For $D=1$ use the trivial character
and the same omitted support. This follows from the same repeated-two
norm identity even when $D=4|h|$. One must do this before using a
Gauss basis; the Gauss sum of an imprimitive character modulo a
larger cyclotomic level can be zero.

Theorem~9.7, p.~189, identifies the dual exponential of $z_D$ with
the retained-smoothed de Rham zeta element. Dual exponential commutes
with corestriction, whose de Rham map is field trace. Thus the
coefficient of $e_\chi$ is the exact unnormalized sum
\[
                    \sum_{g\in\operatorname{Gal}(\mathbf Q(\zeta_D)/\mathbf Q)}
                       \chi(g)\,g(\exp^*z_D).             \tag{K12}
\]
Theorem~6.6, p.~163, evaluates precisely this sum. In weight two,
$r=r'=1$, its Betti sign is $\epsilon=\chi(-1)$ and its power of
$2\pi i$ is zero. If $e_B^\epsilon$ is a fixed rational generator
of that Betti line and
$\delta(f,1,\xi)^\epsilon=b_\epsilon e_B^\epsilon$, then
$b_\epsilon\in\mathbf Q^\times$ is fixed and
\[
     \operatorname{per}_f\bigl((\text{K12})\bigr)^\epsilon
          =s_F L_{\operatorname{prime}(2Nd)}(f,\chi,1)\,
                        b_\epsilon e_B^\epsilon.          \tag{K13}
\]
Choose the original de Rham generator so that its Betti period
on this line is $\Omega_E^\epsilon e_B^\epsilon$, absorbing a
fixed rational comparison in $b_\epsilon$. Formula (K13) says
that the scalar multiplying $e_\chi$ in (K12) is
$s_F b_\epsilon L_{\operatorname{prime}(2Nd)}(f,\chi,1)/
\Omega_E^\epsilon$. Its Galois character is $\chi$.

To descend it, take the invariant basis $\sqrt h\,e_\chi$ of the
quadratic character line. With the usual complex convention
$\Omega_E^-=i|\Omega_E^-|$ and $\sqrt h=i\sqrt{|h|}$ for $h<0$,
division by $\sqrt h$ gives, up to the already fixed rational
period choice, exactly
\[
             s_F b_\epsilon\,\frac{L_*(h)}{h}
       \prod_{p\mid 2N}\text{(the omitted local factor at $p$)}.
                                                               \tag{K14}
\]
The signed integer $h$ is odd, so $h^{-1}$ has valuation zero.
This computation also explains the usual Gauss-basis description:
for the primitive character, $\tau(\chi)^2=\Delta_h$ and
$\tau(\chi)e_\chi$ is invariant; replacing it by $\sqrt h e_\chi$
costs the fixed dyadic factor $\sqrt{\Delta_h/h}$. It does not cost
$\varphi(D)$ or any power of the number of prime divisors of $h$.
The discriminant-$4h$ case has ratio two in this last basis change,
and the direct basis calculation (K14) already includes it.
One can verify the normalization without choosing an idempotent.
Write $\beta=\exp^*z_D$ in the original de Rham line and put
$\beta_\chi=\sum_g\chi(g)g\beta$. Then
$u\beta_\chi=\chi(u)\beta_\chi$ for every $u$ in the finite
cyclotomic group. Thus $\beta_\chi\otimes e_\chi$ is invariant,
and its coordinate in the invariant basis
$\sqrt h\,e_\chi$ is $\beta_\chi/\sqrt h$.
Theorem~6.6 uses precisely $\beta_\chi$, rather than
$\beta_\chi/\varphi(D)$. Theorem~9.7 applies here since
$1\le r=r'=1\le k-1=1$; it has no restriction that $p$ be odd.
The dual form $f^*$ in Theorem~6.6 equals $f$, because the elliptic
newform has rational coefficients.

In the real-period convention,
$L_*(h)=\sqrt{|h|}L(E^h,1)/|\Omega_E^\epsilon|$, up to the
fixed rational period comparison already included in $b_\epsilon$.
For $h>0$,
$\Omega_E^+\sqrt h=|\Omega_E^+|\sqrt{|h|}$; for $h<0$,
$\Omega_E^-\sqrt h=-|\Omega_E^-|\sqrt{|h|}$.
Consequently
\[
 \frac{L(E^h,1)}{\Omega_E^\epsilon\sqrt h}
     =\frac{L_*(h)}{h}
\]
with the same fixed period comparison. This proves the signed
factor in (K14), not only its valuation. For a primitive quadratic
character, $\tau(\chi)^2=\Delta_h$ and
$\tau(\chi)/\sqrt h$ equals a fixed choice of
$\sqrt{\Delta_h/h}$. Changing from its Gauss basis to the displayed
invariant basis multiplies the coordinate by this ratio, which
is one or two up to sign. In particular the discriminant-$4h$
case adds no factor depending on the number of odd conductor primes.


